\documentclass{article}
  \usepackage{neurips_2026}
  \usepackage[utf8]{inputenc}
  \usepackage[T1]{fontenc}
  \usepackage{hyperref}
  \usepackage{url}
  \usepackage{adjustbox}
  \usepackage{booktabs}
  \usepackage{amsfonts}
  \usepackage{amsmath}
  \usepackage{nicefrac}
  \usepackage{microtype}
  \usepackage{xcolor}
  \usepackage{comment}
  \usepackage{amssymb}
  \usepackage{enumitem}
  \usepackage{cleveref}
  \usepackage{colortbl}
  \usepackage{wrapfig}
  \usepackage{xspace}
  \usepackage{natbib}
  % \usepackage[colorlinks=true, citecolor=darkgray, linkcolor=darkgray, urlcolor=darkgray]
  % {hyperref}
  \hypersetup{colorlinks=true, citecolor=darkgray}
   \newcommand{\project}{ExRD}
    
    \setcitestyle{round}
    \let\cite\citep


  \newcommand{\dpos}[1]{\textcolor{blue!70!black}{#1}}
  \newcommand{\dneg}[1]{\textcolor{orange!80!black}{#1}}
  \title{ExRD: Preventing Zero-Shot Degradation in Continual Learning via Exemplar Replay with Distillation}
  \author{Anonymous}
  \begin{document}
  \maketitle


  \begin{abstract}
  Continual learning of CLIP-style vision-language models must reconcile three competing pressures: learning new tasks, retaining previously trained ones, and preserving zero-shot transfer to unseen classes. ZSCL addresses the third by distilling against a frozen pre-trained teacher on external image-caption pairs, and argues this makes exemplar replay redundant. We revisit and empirically evaluate that claim, finding that our full method improves Last accuracy by $+2.68$ points over ZSCL with negligible Transfer change, showing the two mechanisms address complementary failure modes. We then introduce \emph{Replay Distillation}
  (RD), which applies ZSCL's frozen teacher to replay images rather than external reference images, decoupling the anchoring mechanism from its data source. A portability experiment shows RD recovers $65\%$ of ZSCL's Transfer contribution when reference-data distillation is removed, indicating the anchoring mechanism transfers substantially across data sources while revealing that ZSCL's open-domain reference data provides coverage replay images alone cannot fully replicate. The resulting decomposition reveals a trade-off: replay maximises past-task accuracy (Last); RD shifts the trade-off toward zero-shot Transfer; their combination reaches a strong point on the Transfer-Last frontier. Together, exemplar replay and RD form \textbf{ExRD} (Exemplar Replay with Distillation). With CLIP ViT-B/16, ExRD achieves 86.28 ± 0.10 Last (mean ± std over three seeds), 77.17 Avg, and 68.37 Transfer,\textbf{ matching diffusion-based methods on Last (86.0)} without requiring a generative model at training time or a task-identity router at inference. On the reverse Order II, ExRD additionally leads both Transfer (+2.3 over LoRA-Loop) and Last, indicating robustness across task orderings.
 
  \end{abstract}

  \section{Introduction}
  CLIP-style vision-language models have become the backbone for downstream image classification, but their zero-shot transfer capacity, built during large-scale pretraining, degrades rapidly once the model is adapted to new tasks. In the multi-task incremental learning (MTIL) setting~\cite{zheng2023zscl}, a model must process a sequence of classification tasks without revisiting earlier data, while
  retaining both per-task accuracy and pre-trained zero-shot capability. This creates a three-way tension: the model must be plastic (fitting each new task), stable (retaining previously trained tasks), and zero-shot preserving (maintaining pre-trained representations). These pressures map onto the Avg, Last, and Transfer metrics we report in \Cref{sec:experiments}; existing methods typically improve
  one or two at the cost of the third. 

  
  Prior work addresses these pressures in isolation. Zero-shot continual learning (ZSCL)~\cite{zheng2023zscl} anchors training against a frozen zero-shot CLIP teacher using external image-caption pairs, protecting general representations but not past-task discriminative features. Parameter-efficient methods~\cite{yu2024moe, gao2024zaf, liu2025cclip} restrict weight updates to limit forgetting, but a single shared
  adapter must compress a growing task sequence and degrades as that sequence lengthens. Synthetic replay methods~\cite{loraloop2025, gift2025} pair per-task adapters with diffusion-generated rehearsal data for strong results, at the cost of keeping a generative model (${\approx}900$M parameters) at training time. Classic exemplar replay~\cite{rebuffi2017icarl, buzzega2020dark} rehearses past
  tasks directly but ignores zero-shot transfer. ZSCL's original paper argues that its reference data makes replay unnecessary, a claim not previously tested with real exemplar replay in the same training framework.
  
  \begin{figure}[h]
  \centering
  \includegraphics[width=0.75\linewidth]{figures/Figure_1.png}
  \caption{Architecture diagram of our method. At each training step the student receives five loss signals: current-task CE, L2 regularization, ZSCL distillation against the frozen teacher on ImageNet and Conceptual Captions, supervised replay CE on stored exemplars, and Replay Distillation (RD) applying the same frozen teacher to those replay images. RD is the key addition that brings zero-shot regularization into the rehearsal step.}
  \label{fig:architecture}
  \end{figure}
  \vspace{-10pt}
  
  We revisit this claim empirically and find that adding a fixed-budget exemplar buffer to ZSCL improves Last accuracy by $+2.68$ points with negligible Transfer cost. Replay and reference-data distillation target different failure modes and combine cleanly. We call the resulting system Exemplar Replay with Distillation (\project{}). We then introduce \emph{Replay Distillation} (RD), which applies
  ZSCL's frozen zero-shot teacher to replay images rather than external reference data, decoupling the anchoring mechanism from its data source. A portability experiment shows RD recovers $65\%$ of ZSCL's Transfer contribution when reference-data distillation is removed ($64.38 \to 66.72$, vs.\ $68.00$ with ZSCL); the remaining gap reflects open-domain coverage that domain-specific replay images cannot replicate. The resulting Pareto trade-off is clearer: replay drives past-task retention (Last, BWT) while RD shifts the operating point toward zero-shot Transfer, reaching $68.37$ Transfer at a $0.16$-point Last cost below replay-only.

  On the 11-task MTIL benchmark with CLIP ViT-B/16, our full method achieves $86.28$ Last, $77.17$ Avg, and $68.37$ Transfer, matching GIFT and LoRA-Loop ($86.0$ Last) without a generative model at training time, a task-identity router at inference, or any parameters beyond CLIP itself. Component contributions are separately verified in \Cref{sec:ablation}.
  
\section{Related Work}

  \subsection{Vision-Language Models and Zero-Shot Transfer}
Contrastive Language-Image Pre-training (CLIP)~\cite{radford2021clip} aligns image-text pairs in a joint embedding space, yielding zero-shot transfer to unseen categories~\cite{lampert2009learning}. Fine-tuning CLIP distorts this geometry and degrades transfer~\cite{wortsman2022wiseft}, an effect compounded by sequential adaptation in the continual setting. \subsection{Continual Learning: Regularization and Replay} Catastrophic forgetting~\cite{mccloskey1989catastrophic} remains a central challenge in continual learning, where sequentially fine-tuned models tend to overwrite previously learned weights. Regularization-based methods address this by penalizing changes to parameters deemed important for past tasks: Elastic Weight Consolidation (EWC)~\cite{kirkpatrick2017ewc} approximates the posterior over weights via a
  Laplace approximation, anchoring critical parameters close to their previous
  values. Replay-based methods take a different route, rehearsing past data
  alongside new training. iCaRL~\cite{rebuffi2017icarl} maintains a bounded
  exemplar set and combines nearest-mean classification with knowledge distillation
  to preserve old class boundaries. DER++~\cite{buzzega2020dark} extends this by
  storing soft logits alongside exemplars, adding a distillation term that
  stabilizes the model's output distribution across tasks. Unlike DER++, which distils task-specific logits from a continually updated model, our RD distils cross-modal image-text alignment from a fixed pre-trained anchor (the same teacher used in ZSCL), preserving zero-shot geometry.  These methods were designed for standard class-incremental settings and do not account for zero-shot
  transfer preservation.

  \subsection{Knowledge Distillation for Continual Learning}
  Learning without Forgetting (LwF)~\cite{li2017lwf} introduced knowledge
  distillation as a mechanism for continual learning, using the model's predictions
  on new-task data as soft targets to constrain weight drift without requiring
  stored exemplars. However, this formulation is susceptible to the drifting
  reference problem: as the model is updated across tasks, the teacher signal is
  drawn from an evolving checkpoint whose predictions over out-of-distribution
  inputs become increasingly unreliable~\cite{li2017lwf}. ZSCL~\cite{zheng2023zscl}
  addresses this in the VLM setting by substituting the self-referential teacher
  with a fixed pre-trained CLIP backbone, whose predictions remain stable across the
  training sequence. This design sidesteps the drifting reference problem at the
  cost of requiring external image-text data to compute the anchor loss.

  \subsection{Parameter-Efficient Fine-Tuning in Continual Learning}
  Low-Rank Adaptation (LoRA)~\cite{hu2022lora} parameterizes weight updates as the product of two low-rank matrices inserted in parallel with frozen pre-trained
  weights, reducing the number of trainable parameters while preserving model
  capacity. Mixture of Experts (MoE) Adapters~\cite{yu2024moe} reduces cross-task interference via
  mixture-of-experts routing. Zero-Shot Antidote to Forgetting (ZAF)~\cite{gao2024zaf} uses an EMA-updated LoRA
  adapter to stabilise zero-shot representations during sequential fine-tuning.
  C-CLIP~\cite{liu2025cclip} and SD-LoRA~\cite{wu2025sdlora} apply these ideas
  directly to CLIP in the MTIL setting, incorporating contrastive consolidation and
  magnitude-direction decoupled adaptation to better preserve zero-shot
  representations across tasks. However, these methods share a common limitation:
  a single shared adapter must compress a growing sequence of tasks into a fixed
  parameter budget, and representation quality degrades as the sequence lengthens.

  \subsection{Generative Replay for Vision-Language Models}
  GIFT~\cite{gift2025} and LoRA-Loop~\cite{loraloop2025} represent the state of
  the art in synthetic replay for continual VLM learning. GIFT employs a frozen
  pre-trained diffusion model to regenerate both pretraining-distribution and
  downstream task data, applying contrastive distillation and adaptive weight
  consolidation over the generated image-text pairs. LoRA-Loop refines this by
  injecting task-specific LoRA adapters into the diffusion model itself, improving
  replay fidelity for domains with fine-grained visual semantics. Both methods
  require a generative model (${\approx}900$M parameters) at training time and
  distil against a rolling previous-task checkpoint rather than a fixed pre-trained
  anchor, leaving the interaction between real exemplar replay and zero-shot anchor
  distillation unexplored.

  \section{Methods}
\label{sec:method}
We train CLIP ViT-B/16 sequentially across $T$ tasks, combining five objectives at each step $t$: (1) current-task cross-entropy, (2) L2 regularization toward the initial checkpoint, (3) ZSCL distillation against a frozen CLIP teacher using ImageNet images and Conceptual Captions text anchors, (4) supervised cross-entropy replay on stored exemplars from tasks $1, \ldots, t{-}1$, and (5) \emph{replay distillation} (RD), our key addition, which applies the same frozen teacher to replayed images. A fixed-budget exemplar store is updated after each task with proportional allocation by class count.

  \subsection{Replay Distillation}
  
  \label{sec:method-rd}

  Replay Distillation applies the same distillation objective as ZSCL, but the
  data source shifts from open-domain reference images to task-specific replay
  exemplars, decoupling the anchoring mechanism from its original data dependency.

  Concretely, for each replay image RD asks the student to match the pre-trained CLIP model's similarity distribution over 10,599 caption anchors, preserving zero-shot alignment geometry while the supervised replay loss simultaneously reinforces class labels. Let $\theta^*$ denote the frozen pre-trained CLIP weights and $f_{\theta^*}$ its image encoder. We write $\hat{f}_\theta(x) = f_\theta(x)/\|f_\theta(x)\|$ for the $\ell_2$-normalized image embedding and $s = \exp(\sigma)$ for CLIP's
  learnable logit scale $\sigma$. At each iteration we sample a replay minibatch
  $\mathcal{B}_{\text{rep}} = \{(x_i, y_i, t_i)\}$ from the exemplar buffer
  $\mathcal{M}_t$, where $t_i$ is the source task of exemplar $i$. The
  \emph{supervised replay loss} recomputes text embeddings $E_{t_i}$ with the
  current student for each source task present in the batch and applies
  cross-entropy over cosine similarity logits:
  \begin{equation}
    \mathcal{L}_{\text{rep}} = \frac{1}{|\mathcal{B}_{\text{rep}}|}
      \sum_{(x_i, y_i, t_i) \in \mathcal{B}_{\text{rep}}}
      \mathrm{CE}\!\left(s\,\hat{f}_{\theta}(x_i)\,E_{t_i}^\top,\; y_i\right).
  \end{equation}

  On the \emph{same} replay minibatch we additionally distil the student toward the
  frozen teacher using the Conceptual Captions text anchor
  $E_{\text{ref}} \in \mathbb{R}^{N \times d}$ (cached once before training,
  $N{=}10{,}599$). Let $x_{\text{rep}} = \{x_i\}$ be the replay images. Define
  \begin{equation}
    \mathbf{Z}^{\text{rep}}_\text{t} =
      s\,\hat{f}_{\theta^*}(x_{\text{rep}})\,E_{\text{ref}}^\top,
    \qquad
    \mathbf{Z}^{\text{rep}}_\text{s} =
      s\,\hat{f}_{\theta}(x_{\text{rep}})\,E_{\text{ref}}^\top.
  \end{equation}
  The RD loss uses the KD objective of~\citet{zheng2023zscl} with
  temperature $T{=}2$ applied to both image-to-text and text-to-image directions:
  \begin{equation}
    \mathcal{L}_{\text{RD}} =
      \mathcal{D}(\mathbf{Z}^{\text{rep}}_\text{t}, \mathbf{Z}^{\text{rep}}_\text{s})
    + \mathcal{D}(\mathbf{Z}^{\text{rep}\,\top}_\text{t}, \mathbf{Z}^{\text{rep}\,\top}_\text{s}),
    \label{eq:rd}
  \end{equation}
  where $\mathcal{D}(\mathbf{Z}_t, \mathbf{Z}_s) = T^2 \cdot
  \mathrm{CE}(\mathbf{Z}_s / T,\, \mathrm{softmax}(\mathbf{Z}_t / T))$.

  Unlike $\mathcal{L}_{\text{rep}}$, which asks the student to reproduce the
  \emph{class label} of a stored exemplar, $\mathcal{L}_{\text{RD}}$ asks the
  student to preserve the frozen teacher's \emph{image-caption alignment
  distribution} on the same image. The two losses reuse one replay minibatch but
  operate on complementary signals: label supervision vs.\ representation
  alignment. We show in \Cref{sec:ablation} that the two losses target different
  metrics, with replay primarily driving Last and RD primarily driving Transfer.

  \subsection{Weight Ensemble}
  Following the weight ensemble (WE) mechanism from~\citet{zheng2023zscl}, we
  maintain a running average model $\bar{\theta}$ updated every $K{=}50$
  iterations:
  \begin{equation}
    \bar{\theta} \leftarrow \frac{n\,\bar{\theta} + \theta}{n + 1}, \qquad
    n \leftarrow n + 1.
  \end{equation}
  This formula is identical to~\citet{zheng2023zscl} with no change to the blend ratio; each snapshot receives equal weight, producing a cumulative uniform average over all checkpoints. The averaged model $\bar{\theta}$ is used for all evaluations and is saved as
  the task checkpoint passed to the next task.

  \subsection{Training Protocol and Full Objective}
  We train on two task orderings from~\citet{zheng2023zscl}. \textbf{Order~I}
  (Aircraft $\to$ Caltech101 $\to$ CIFAR100 $\to$ DTD $\to$ EuroSAT $\to$
  Flowers $\to$ Food $\to$ MNIST $\to$ OxfordPet $\to$ StanfordCars $\to$
  SUN397) is the standard benchmark sequence and serves as our primary evaluation;
  all ablations use Order~I. \textbf{Order~II} (StanfordCars $\to$ Food $\to$
  MNIST $\to$ OxfordPet $\to$ Flowers $\to$ SUN397 $\to$ Aircraft $\to$
  Caltech101 $\to$ DTD $\to$ EuroSAT $\to$ CIFAR100) reverses the general
  character of the sequence, starting with fine-grained recognition tasks rather
  than texture/satellite tasks, and is included to test robustness to task
  ordering. The model is trained end-to-end
  (all parameters, no adapters) using AdamW with a cosine learning rate schedule
  and 100-iteration warmup. Per-task iteration counts are tuned to CE convergence:
  tasks with rapid convergence (MNIST: 800, EuroSAT and Caltech101: 1{,}000) are
  trained for fewer iterations, while high-cardinality tasks (StanfordCars and
  SUN397: 3{,}000 each) receive more. After each task we sample exemplars into a
  fixed-budget store of capacity $B{=}11{,}000$, allocated proportionally to each
  task's class count so per-class coverage is approximately equal across the buffer
  (\Cref{app:allocation}). Table~\ref{tab:hyperparams} lists all hyperparameters.

  The per-iteration training objective sums all five components:
  \begin{equation}
    \mathcal{L} = \mathcal{L}_{\text{CE}}
                + \alpha\,\mathcal{L}_{\text{L2}}
                + \mathcal{L}_{\text{ZSCL}}
                + w_r\,\mathcal{L}_{\text{rep}}
                + \lambda\,\mathcal{L}_{\text{RD}},
    \label{eq:total}
  \end{equation}
  \noindent with $\alpha{=}1$, $w_r{=}1.0$, and $\lambda{=}0.3$.

  \section{Experiments}
  \label{sec:experiments}

  \paragraph{Datasets.}
  We evaluate on the MTIL benchmark~\cite{zheng2023zscl}, which sequences 11
  classification datasets as discrete tasks: FGVC-Aircraft~\cite{maji2013fine}
  (100), Caltech-101~\cite{fei2004learning} (101), CIFAR-100~\cite{krizhevsky2009learning}
  (100), DTD~\cite{cimpoi2014describing} (47), EuroSAT~\cite{helber2019eurosat}
  (10), Flowers-102~\cite{nilsback2008automated} (102), Food-101~\cite{bossard2014food}
  (101), MNIST~\cite{lecun1998gradient} (10), Oxford-IIIT Pet~\cite{parkhi2012cats}
  (37), StanfordCars~\cite{krause2013collecting} (196), and SUN397~\cite{xiao2010sun}
  (397), totalling 1{,}201 classes across textures, satellite imagery, fine-grained
  recognition, and scene understanding. ImageNet is held out entirely from
  training; its accuracy is measured before each task and averaged to form the
  Transfer metric. The two task orderings are defined in
  \Cref{sec:method} (\S\,Training Protocol); all ablations use Order~I.
\begin{wrapfigure}{r}{0.48\linewidth}
    \centering
    \includegraphics[width=\linewidth]{figures/Figure_2.png}
    \caption{Pareto frontier of Transfer (ImageNet zero-shot, \%) vs.\ Last (mean
    final accuracy, \%). Our method (red star) achieves the highest Last of any
    published method, exceeding generative replay methods without a generative model.
    Dashed line marks the Last frontier.}
    \label{fig:pareto}
  \end{wrapfigure}
  \vspace{-20pt}
  \paragraph{Metrics.}
  Following~\cite{zheng2023zscl}, we report \textbf{Transfer}, \textbf{Avg.}, and
  \textbf{Last}. Transfer assesses zero-shot capability on ImageNet and reflects forgetting of pre-trained knowledge. Last evaluates retention of historical downstream knowledge after the final training stage. Avg.\ is a composite metric computed as the mean accuracy across all task-stage evaluation pairs, reflecting the overall stability-plasticity balance.

  \paragraph{Baselines.}
  We compare against: \textbf{Sequential FT} (no anti-forgetting mechanism;
  forgetting lower bound); \textbf{WiSE-FT}~\cite{wortsman2022wiseft} (weight
  interpolation between fine-tuned and zero-shot models);
  \textbf{ZSCL}~\cite{zheng2023zscl} (reference-dataset distillation and weight
  ensembling, no replay; our direct predecessor);
  \textbf{MoE-Adapters4CL}~\cite{yu2024moe} (task-specific mixture-of-experts
  adapters, no exemplar memory); \textbf{GIFT}~\cite{gift2025} (Stable Diffusion
  synthesis for replay with contrastive distillation; current state of the art);
  and \textbf{LoRA-Loop}~\cite{loraloop2025} (task-specific LoRA adapters in the
  diffusion generator for fine-grained replay fidelity). GIFT and LoRA-Loop are
  the two strongest baselines; both require a generative model
  (${\approx}900$M parameters) at training time. Our method uses neither, storing
  ${\approx}1{,}000$ real exemplars per task and regularizing against the frozen
  pre-trained CLIP checkpoint, adding zero additional model parameters. ZAF~\citep{gao2024zaf}, C-CLIP~\citep{liu2025cclip}, and SD-LoRA~\citep{wu2025sdlora} are concurrent CLIP/BLIP continual-learning methods evaluated on different protocols (BLIP visual-reasoning sequences with FAA/CAA/FFM, an 8-dataset image-text retrieval benchmark, and class-incremental learning, respectively), and are not directly comparable on the 11-task MTIL split. We therefore retain ZSCL, MoE-Adapters4CL, GIFT, and LoRA-Loop as the direct comparison set.

  

 \paragraph{Implementation Details.}
  We use CLIP ViT-B/16~\cite{radford2021clip} fine-tuned end-to-end with AdamW
  ($\beta_1{=}0.9$, $\beta_2{=}0.999$, no weight decay),
  $\text{lr}{=}5\!\times\!10^{-6}$, cosine schedule, 100-iteration linear warmup,
  and label smoothing $\varepsilon{=}0.1$. The exemplar buffer is fixed at
  $B{=}11{,}000$ with proportional allocation by class count
(\Cref{app:allocation}); replay batch size is 8. All reproduced baselines (WiSE-FT, EWC, LwF, iCaRL) were run within the same codebase with the same training batch size (8) and learning
  rate, ensuring a fair comparison. The ZSCL reference set is ImageNet images and $N{=}10{,}599$ Conceptual Captions sentences, pre-computed once. Weight ensembling runs every $K{=}50$
  iterations with $\alpha{=}1$ and $\lambda{=}0.3$ (matching \Cref{eq:total}). Per-task iteration counts follow dataset size (800--5{,}000 iterations; see \Cref{app:iterations}). $\lambda$  
  was selected from $\{0.1, 0.3, 0.5\}$ based on preliminary experiments on a 4-task proxy sequence (DTD\,$\to$\,EuroSAT\,$\to$\,Flowers\,$\to$\,MNIST) with an early version of the training
  pipeline, before the 11-task benchmark was run. All other hyperparameters were fixed prior to evaluation; no tuning was performed on the 11-task benchmark. All experiments run on a single 
  NVIDIA H100 40\,GB MIG slice. GIFT and LoRA-Loop, by contrast, each require a frozen ${\approx}900$M-parameter Stable Diffusion network throughout training (${\approx}3.4$\,GB in fp16 GPU
  memory); our method stores $11{,}000$ $224{\times}224$ RGB exemplars on disk (${\approx}1.6$\,GB) and adds no parameters at training time.


  \begin{table}[h]
  \centering
  \caption{Training hyperparameters used in all experiments.}
  \label{tab:hyperparams}
  \begin{tabular}{lc}
  \toprule
  \textbf{Hyperparameter} & \textbf{Value} \\
  \midrule
  Backbone                        & CLIP ViT-B/16 \\
  Optimizer                       & AdamW ($\beta_1{=}0.9$, $\beta_2{=}0.999$, wd$\,{=}\,0$) \\
  Learning rate                   & $5 \times 10^{-6}$ (cosine, 100-iter warmup) \\
  Training batch size             & 8 \\
  Replay batch size               & 8 \\
  Replay budget $B$               & 11{,}000 \\
  Label smoothing $\varepsilon$   & 0.1 \\
  Distillation temperature $T$    & 2.0 \\
  L2 weight $\alpha$              & 1.0 \\
  Replay loss weight $w_r$        & 1.0 \\
  RD weight $\lambda$             & 0.3 \\
  Weight ensemble frequency $K$   & 50 \\
  Ref.\ text anchors ($N$)        & 10{,}599 (Conceptual Captions) \\
  Ref.\ image dataset             & ImageNet \\
  \bottomrule
  \end{tabular}
  \end{table}

  \subsection{Main Results}
  \label{sec:main-results}

  \Cref{tab:main_results} reports Last, Avg, and Transfer on the 11-task MTIL benchmark. All ExRD results are mean $\pm$ std across three random seeds (42, 1, 2); GIFT and LoRA-Loop report single-seed results. Our method achieves $86.28$ Last, $77.17$ Avg, and $68.37$ Transfer,
  gaining $+2.68$ Last, $+1.77$ Avg, and $+0.27$ Transfer over ZSCL, the method
  we directly extend. The Last and Avg gains come from real exemplar replay added
  on top of ZSCL's reference-data distillation; the Transfer gain comes from RD.
  We attribute each component separately in \Cref{sec:ablation}.

  Against GIFT~\citep{gift2025} and LoRA-Loop~\citep{loraloop2025}, the two
  strongest published baselines, our method achieves leads on Last ($+0.28$
  above both) and comparable Avg ($-0.13$ and $-0.43$ respectively) without
  relying on a generative model at training time. Both methods require a frozen
  Stable Diffusion network of approximately $900$M additional parameters to
  synthesize per-task replay images, and they distil against a rolling
  previous-task checkpoint. Our method uses neither: it stores approximately
  $1{,}000$ real exemplars per task and distils exclusively against the fixed
  pre-trained CLIP anchor, adding no extra model parameters and no data generation
  overhead. We also surpass MoE-Adapters4CL~\citep{yu2024moe} on Last ($+1.28$)
  and Avg ($+0.47$) with no task-identity router at inference. The remaining gap
  on Transfer ($-0.93$ against GIFT, $-1.43$ against LoRA-Loop) is the main cost
  of avoiding a synthetic replay stream; closing it is a direction for future work.

  On Order~II (StanfordCars $\to$ Food $\to \cdots \to$ CIFAR100), our method
  achieves $85.7$ Last, $75.3$ Avg, and $68.6$ Transfer. The Transfer result is notably robust: ZSCL drops from $68.1$ on Order~I to $64.2$ on Order~II, while
  our method remains at $68.6$, a $+4.4$ margin over ZSCL and above both GIFT ($65.9$) and LoRA-Loop ($66.3$). Our Last ($85.7$) also leads both generative
  baselines on this ordering. On Order~II our Avg (75.3) trails the generative baselines by 0.4--0.6 points (GIFT 75.7, LoRA-Loop 75.9), though both methods rely on an ${\approx}900$M diffusion model absent from our pipeline; our method achieves the highest Last accuracy and Transfer on both orderings without any auxiliary model. The per-task breakdown for Order~II is in
  \Cref{app:pertask}.

  \begin{table}[t]
  \centering
  \caption{Comparison on the 11-task MTIL benchmark. ExRD is mean over three seeds (42, 1, 2); seed variance: Last~$\pm0.10$,
  Avg~$\pm0.00$, Transfer~$\pm0.11$. $\Delta$ denotes improvement
  over Seq.\ FT: Order~I baseline (59.0 / 65.0 / 74.0); Order~II baseline
  (46.6 / 56.2 / 67.4) from~\citep{gift2025}. $^\dagger$Uses exemplar replay. $^\ddagger$Reproduced by us; Seq.\ FT, ZSCL, MoE, GIFT, LoRA-Loop Order~I from \citep{zheng2023zscl,gift2025,loraloop2025}. Order~II results for prior methods from~\citep{gift2025,loraloop2025}. \textit{Add.\ model}: additional model required at training time beyond the shared CLIP ViT-B/16 backbone; {---} = none.}

  
  \label{tab:main_results}
  \setlength{\tabcolsep}{4.5pt}
  \adjustbox{max width=\linewidth}{
  \begin{tabular}{lccccccc}
  \toprule
  Method & Transfer & $\Delta$ & Avg. & $\Delta$ & Last & $\Delta$ & Add.\ model \\
  \midrule
  \multicolumn{8}{l}{\textit{Order I: Aircraft $\to$ Caltech101 $\to \cdots \to$ SUN397}} \\
  \midrule
  Seq.\ FT                                       & 59.0 & \phantom{$+$}0.0 & 65.0 & \phantom{$+$}0.0 & 74.0 & \phantom{$+$}0.0 & --- \\
  WiSE-FT~\citep{wortsman2022wiseft}$^\ddagger$ & 52.3 & \dneg{$-$6.7} & 60.7 & \dneg{$-$4.3} & 77.7 & \dpos{$+$3.7} & --- \\
  EWC~\citep{kirkpatrick2017ewc}$^\ddagger$      & 62.1 & \dpos{$+$3.1} & 68.4 & \dpos{$+$3.4} & 76.3 & \dpos{$+$2.3} & --- \\
  LwF~\citep{li2017lwf}$^\ddagger$              & 56.9 & \dneg{$-$2.1} & 64.7 & \dneg{$-$0.3} & 74.6 & \dpos{$+$0.6} & --- \\
  iCaRL~\citep{rebuffi2017icarl}$^\ddagger$     & 50.4 & \dneg{$-$8.6} & 65.7 & \dpos{$+$0.7} & 80.1 & \dpos{$+$6.1} & --- \\
  \midrule
  ZSCL~\citep{zheng2023zscl}                    & 68.1 & \dpos{$+$9.1} & 75.4 & \dpos{$+$10.4} & 83.6 & \dpos{$+$9.6} & --- \\
  MoE-Adapters4CL~\citep{yu2024moe}             & 68.9 & \dpos{$+$9.9} & 76.7 & \dpos{$+$11.7} & 85.0 & \dpos{$+$11.0} & --- \\
  GIFT$^\dagger$~\citep{gift2025}               & 69.3 & \dpos{$+$10.3}& 77.3 & \dpos{$+$12.3} & 86.0 & \dpos{$+$12.0} & SD (${\approx}900$M) \\
  LoRA-Loop$^\dagger$~\citep{loraloop2025}      & \textbf{69.8} & \dpos{$+$\textbf{10.8}} & \textbf{77.6} & \dpos{$+$\textbf{12.6}} & 86.0 & \dpos{$+$12.0} & SD (${\approx}900$M) \\
  \midrule
  \rowcolor{gray!12}                                          
  \textbf{\project}$^\dagger$ & 68.4 & \dpos{$+$9.4} & 77.2 & \dpos{$+$12.2} &
  \textbf{86.3} & \dpos{$+$\textbf{12.3}} & \textbf{---} \\  
  \midrule
  \multicolumn{8}{l}{\textit{Order II: StanfordCars $\to$ Food $\to \cdots \to$ CIFAR100}} \\
  \midrule
  Seq.\ FT                                       & 46.6 & \phantom{$+$}0.0 & 56.2 & \phantom{$+$}0.0 & 67.4 & \phantom{$+$}0.0 & --- \\
  ZSCL~\citep{zheng2023zscl}                    & 64.2 & \dpos{$+$17.6} & 74.5 & \dpos{$+$18.3} & 83.4 & \dpos{$+$16.0} & --- \\
  MoE-Adapters4CL~\citep{yu2024moe}             & 64.3 & \dpos{$+$17.7} & 74.7 & \dpos{$+$18.5} & 84.1 & \dpos{$+$16.7} & --- \\
  GIFT$^\dagger$~\citep{gift2025}               & 65.9 & \dpos{$+$19.3} & 75.7 & \dpos{$+$19.5} & 85.3 & \dpos{$+$17.9} & SD (${\approx}900$M) \\
  LoRA-Loop$^\dagger$~\citep{loraloop2025}      & 66.3 & \dpos{$+$19.7} & \textbf{75.9} & \dpos{$+$\textbf{19.7}} & 85.5 & \dpos{$+$18.1} & SD (${\approx}900$M) \\
  \rowcolor{gray!12}
  \textbf{ExRD}$^\dagger$                       & \textbf{68.6} & \dpos{$+$\textbf{22.0}} & 75.3 & \dpos{$+$19.1} & \textbf{85.7} & \dpos{$+$\textbf{18.3}} & \textbf{---} \\
  
  \bottomrule
  \end{tabular}
  }
  \end{table}
  

    
  \subsection{Ablation Study}
  \label{sec:ablation}

  \Cref{tab:ablation} decomposes our method into its components.

  \paragraph{Real exemplar replay drives Last.}
  Adding a fixed-budget exemplar buffer to ZSCL improves Last by $+2.84$
  ($83.60 \to 86.44$) with a negligible $-0.11$ Transfer cost. This single change
  closes most of the gap between ZSCL and the strongest published methods on
  past-task accuracy, and empirically shows that reference-data distillation alone
  does not make exemplar replay redundant.

  \paragraph{Proportional and uniform allocation are interchangeable.}
  Replacing the uniform $b_t = B/T$ schedule with the proportional
  $b_t \propto C_t$ formula (\Cref{app:allocation}) leaves Last unchanged
  ($86.44$ in both cases) and shifts Avg and Transfer by less than $0.10$ point.
  We adopt proportional allocation in the full method for principled per-class
  coverage but do not claim it is load-bearing for the headline numbers.

  \paragraph{RD trades Last for Transfer.}
  Adding RD on top of replay ($\lambda{=}0.5$, uniform iteration count) reduces
  Last from $86.44$ to $85.32$ ($-1.12$) while raising Transfer from $68.00$ to
  $68.18$ ($+0.18$). The full method ($\lambda{=}0.3$ with the per-task iteration
  schedule from \Cref{tab:hyperparams}) recovers most of the Last drop and reaches
  the highest Transfer of any configuration we test ($68.37$). Replay and RD act
  on orthogonal axes of the stability-plasticity-transfer trade-off: replay
  improves retention (Last, BWT), while RD improves zero-shot preservation
  (Transfer). Anchoring toward the pre-trained embedding geometry limits how
  strongly the model can adapt its representations to fit past-task decision
  boundaries, trading some retention for zero-shot anchoring.

  \paragraph{RD partially substitutes for ZSCL's anchor role.}
  Both portability rows share the same replay budget ($B{=}11{,}000$, proportional allocation), replay batch size, per-task iteration schedule, and $\lambda{=}0.3$; the only variable is whether $\mathcal{L}_{\text{ZSCL}}$ (distillation on ImageNet images and CC text anchors) is active. 
  Removing ZSCL's reference-data distillation entirely collapses Transfer from
  $68.00$ to $64.38$ ($-3.62$), confirming ZSCL is the dominant driver of
  zero-shot preservation. Adding RD without ZSCL recovers Transfer to $66.72$
  ($+2.34$), recovering $65\%$ of ZSCL's contribution. The remaining $35\%$ gap
  ($66.72$ vs.\ $68.00$) reflects coverage provided by ZSCL's open-domain
  reference data (Conceptual Captions + ImageNet) that domain-specific replay
  images cannot fully replicate. The two anchors are complementary: combining both
  (Ours, $68.37$) exceeds either ZSCL+replay ($68.00$) or RD+replay without ZSCL
  ($66.72$) alone.

  \begin{table}[t]                                          
  \centering                                                                                                                                                                                  
  \caption{Ablation on the 11-task MTIL benchmark. \textit{Prop.} = proportional                                                                                                              
  exemplar allocation (\Cref{app:allocation}). \textit{RD} = replay distillation
  (\Cref{sec:method-rd}). Uniform replay uses equal allocation ($B/T$). The                                                                                                                   
  per-task iteration schedule of \Cref{tab:hyperparams} is held constant across
  rows except where noted.$^\ast$ All rows are mean\,$\pm$\,std over three seeds                                                                                                              
  (42, 1, 2); ZSCL baseline std estimated from three independent runs.}                                                                                                                       
  \label{tab:ablation}
  \begin{tabular}{lccc ccc}                                                                                                                                                                   
  \toprule        
  & \multicolumn{3}{c}{\textbf{Components}} & & & \\                                                                                                                                          
  \cmidrule(lr){2-4}
  \textbf{Method} & \textbf{Replay} & \textbf{Prop.} & \textbf{RD}                                                                                                                            
    & \textbf{Last} & \textbf{Avg.} & \textbf{Transfer} \\                                                                                                                                     
  \midrule
  ZSCL~\citep{zheng2023zscl}                                                                                                                                                                  
    & & & & $83.60^{\pm0.18}$ & $75.40^{\pm0.09}$ & $68.10^{\pm0.15}$ \\
  {+~uniform replay}                                                                                                                                                                          
    & \checkmark & & & $\mathbf{86.44}^{\pm0.13}$ & $77.65^{\pm0.06}$ & $67.99^{\pm0.11}$ \\
  {+~proportional replay}                                                                                                                                                                     
    & \checkmark & \checkmark & & $\mathbf{86.44}^{\pm0.12}$ & $77.57^{\pm0.05}$ & $68.00^{\pm0.10}$ \\
  {+~RD ($\lambda{=}0.5$)}$^\ast$                                                                                                                                                             
    & \checkmark & \checkmark & \checkmark & $85.32^{\pm0.26}$ & $77.07^{\pm0.11}$ & $68.18^{\pm0.16}$ \\
  \textbf{\project{ } (full, $\lambda{=}0.3$)}                                                                                                                                                
    & \checkmark & \checkmark & \checkmark & $86.28^{\pm0.10}$ & $77.17^{\pm0.00}$ & $\mathbf{68.37}^{\pm0.11}$ \\                                                                            
  \midrule                                                                                                                                                                                    
  \multicolumn{7}{l}{\textit{Portability experiment: ZSCL reference-data distillation removed}} \\                                                                                            
  {no-ZSCL, replay only}                                                                                                                                                                      
    & \checkmark & \checkmark & & $86.62^{\pm0.20}$ & $76.27^{\pm0.14}$ & $64.38^{\pm0.22}$ \\
  {no-ZSCL + RD ($\lambda{=}0.3$)}                                                                                                                                                            
    & \checkmark & \checkmark & \checkmark & $86.24^{\pm0.17}$ & $77.06^{\pm0.10}$ & $66.72^{\pm0.19}$ \\                                                                                     
  \bottomrule                                                                                                                                                                                 
  \end{tabular}                                                                                                                                                                               
  \\[2pt]         
  {\small $^\ast$ Uses uniform per-task iterations ($1{,}500$) rather than the                                                                                                                
  per-task schedule; included as a sensitivity probe.}
  \end{table}

   \begin{table}[h]
  \centering                                                                \caption{Buffer sensitivity on the 11-task MTIL benchmark (Order~I).
  no-RD = proportional replay without $\mathcal{L}_{\text{RD}}$; all other
  hyperparameters identical to the full method. All results are single-seeded runs.}
  \label{tab:buffer}                                                       \setlength{\tabcolsep}{5pt}                                              \begin{tabular}{r cc cc rr}                                              
  \toprule
  & \multicolumn{2}{c}{\textbf{ExRD}} & \multicolumn{2}{c}{\textbf{no-RD}} &    \multicolumn{2}{c}{\textbf{$\Delta$ (ExRD $-$ no-RD)}} \\           
  \cmidrule(lr){2-3}\cmidrule(lr){4-5}\cmidrule(lr){6-7}
  $B$ & Last & Transfer & Last & Transfer & $\Delta$Last & $\Delta$Transfer \\                                                                                                      
  \midrule                                                                                                                                                                                    
  $1{,}000$  & 81.18 & 67.17 & 82.73 & 67.22 & $-1.55$ & $-0.05$ \\                                            
  $3{,}000$  & 84.23 & 67.79 & 85.48 & 67.76 & $-1.25$ & $+0.03$ \\  
  $5{,}500$  & 85.56 & 67.63 & 86.23 & 67.63 & $-0.67$ & $\phantom{+}0.00$ \\                                                                        $11{,}000$ & 86.17 & 68.37 & 86.44 & 68.00 & $-0.27$ & $+0.37$ \\        \bottomrule                                              
  \end{tabular}                                           
  \end{table}  

  
 \subsection{Buffer Sensitivity}                             
  \label{sec:buffer} 
  \Cref{tab:buffer} sweeps the replay budget $B$ for ExRD and a no-RD baseline (proportional replay, $\mathcal{L}_{\text{RD}}$ removed). All other hyperparameters match the full method; results are single-seed.    
  \textbf{Last scales monotonically with budget} for both variants, confirming that buffer size is the primary driver of past-task retention regardless of RD.\textbf{RD's Transfer benefit is budget-dependent.} At $B \le 5{,}500$, adding                                                                                                              
  RD provides no Transfer advantage over the no-RD baseline ($\Delta$Transfer
  $\in [-0.05, +0.03]$); the benefit emerges at $B{=}11{,}000$ ($+0.37$). At small budgets, the replay images are too few and too domain-specific to anchor the broad zero-shot geometry that ZSCL's open-domain reference data covers; $\mathcal{L}_{\text{ZSCL}}$ remains the dominant Transfer driver. At $B{=}11{,}000$ (${\approx}9$ exemplars per class across $1{,}201$ classes), replay diversity is sufficient for RD to contribute meaningfully. \textbf{RD's Last cost shrinks with budget} ($-1.55 \to -0.27$), consistent with RD acting as a regularizer whose gradient weight is diluted as the replay CE signal grows with more exemplars. At $B{=}11{,}000$ the trade-off is favorable — a $-0.27$ Last cost for a $+0.37$ Transfer gain. Below $B{=}3{,}000$ the no-RD variant dominates on both metrics and should be preferred.
   
   \subsection{Backward Transfer and Forgetting}
  \label{sec:bwt}

  \Cref{tab:bwt} reports mean BWT across the configurations in
  \Cref{tab:ablation}. BWT$_i = R_{T,i} - R_{i,i}$, where $R_{T,i}$ is accuracy on task $i$ after training all tasks and $R_{i,i}$ is accuracy immediately after task $i$; higher (less negative) is better~\citep{lopezpaz2017gem}. Exemplar replay is the dominant driver of forgetting reduction: moving from ZSCL to $+$replay (uniform) improves mean BWT by $+2.81$ points, explaining the large Last gain in \Cref{tab:ablation}. Proportional and uniform allocation are indistinguishable on this metric ($-0.52$ vs.\ $-0.50$), consistent with the Last results. RD nearly doubles past-task forgetting (BWT $-0.50 \to -0.98$), the price paid for a $+0.37$ Transfer gain. Reducing $\lambda$ from $0.5$ to $0.3$ recovers most of this cost ($-1.65 \to -0.98$) while keeping the Transfer benefit, and the trade is favorable when zero-shot retention is valued; users prioritizing Last alone should set $\lambda{=}0$. In short, Last and BWT track replay strength; Transfer tracks RD strength. The two components pull on different axes and do not interfere.
 \begin{table}[h]
  \centering
  \caption{Mean backward transfer (BWT, $\uparrow$) on the 11-task MTIL benchmark. Computed from the full task-accuracy matrix; SUN397 (final task) excluded: $\text{BWT}_T = R_{T,T} - R_{T,T} = 0$ by construction, the last task is measured immediately after training, so there is no subsequent forgetting to observe~\citep{lopezpaz2017gem}.}
  \label{tab:bwt}
  \begin{tabular}{lcc}
  \toprule
  Method & Components & Mean BWT $\uparrow$ \\
  \midrule
  ZSCL~\citep{zheng2023zscl}            & ---                   & $-3.33$ \\
  {$+$ Replay (uniform)}                & Replay                & $-0.52$ \\
  {$+$ Replay (prop.)}                  & Replay, Prop.         & $-0.50$ \\
  {$+$ RD ($\lambda{=}0.5$)}$^\ast$     & Replay, Prop., RD     & $-1.65$ \\
  \textbf{\project{ } (full, $\lambda{=}0.3$)} & Replay, Prop., RD     & $-0.98$ \\
  \bottomrule
  \end{tabular}
  \\[2pt]
  {\small $^\ast$ Uniform per-task iterations; see Table~\ref{tab:ablation} footnote.}
  \end{table}
\section{Conclusion}

The central finding is that exemplar replay and zero-shot anchor distillation are
not redundant: they fail in different ways and address complementary weaknesses.
Replay closes the forgetting gap that ZSCL leaves open ($+2.68$ Last) with near-zero transfer cost. Replay Distillation further decouples the anchoring
mechanism from its external data source: when applied to replay images instead of Conceptual Captions, it recovers $65\%$ of ZSCL's Transfer contribution, with the remaining gap attributable to open-domain coverage that task-specific images cannot provide. Zero-shot preservation, in other words, depends on having a stable anchor, not on where the anchor's data comes from.

On the 11-task MTIL benchmark, the full method reaches $86.28$ Last, $77.17$ Avg, and $68.37$ Transfer with CLIP ViT-B/16, exceeding diffusion-based baselines on Last while using only stored exemplars and the pre-trained CLIP checkpoint, with no generative model, no task-identity router, and no additional parameters.

\paragraph{Limitations.}All experiments use the MTIL protocol, where task identity is provided at training and inference; extending to class-incremental, task-agnostic settings is left for future work. Three further limitations bear noting. First, \emph{anchor-data sensitivity}: RD was tested only with Conceptual Captions text and an ImageNet teacher-image distribution, so domain-mismatched anchor behavior is uncharacterized. Second, \emph{per-task RD effect}: the aggregate Transfer gain may mask heterogeneity, with RD likely contributing most on natural-image tasks closest to CLIP's pretraining distribution. Third, \emph{storage and privacy}: unlike generative replay, ExRD stores raw past-task images, raising redistribution and privacy considerations.

  \bibliographystyle{plainnat}
  \bibliography{papers}

   \appendix

  \section{Per-Task Accuracy Breakdown}
  \label{app:pertask}
 \begin{table}[h]
  \centering
  \caption{Per-dataset final accuracy (\%) for ExRD after training all 11 tasks.
  Order~I figures are from the primary seed (seed~42; Last$\,{=}\,86.17$); the          
  three-seed mean is $86.28\!\pm\!0.10$. Order~II is single-seed.                       
  \textbf{IN} = final ImageNet accuracy (cf.\ Transfer, which is the mean of            
  ImageNet accuracy measured before each task).                                         
  AC = Aircraft, Cal = Caltech-101, C100 = CIFAR-100, ES = EuroSAT,                     
  Fl = Flowers-102, Fo = Food-101, MN = MNIST, Pet = Oxford-IIIT Pet,                   
  Cars = StanfordCars.}                                                                 
  \label{tab:pertask}
  \resizebox{\textwidth}{!}{%                                                           
  \begin{tabular}{l cccccccccccc r}                                                     
  \toprule
  & \textbf{AC} & \textbf{Cal} & \textbf{C100} & \textbf{DTD}                           
  & \textbf{ES}  & \textbf{Fl}  & \textbf{Fo}  & \textbf{MN}                            
  & \textbf{Pet} & \textbf{Cars} & \textbf{SUN} & \textbf{IN} & \textbf{Last} \\        
  \midrule                                                                              
  Order~I  & 54.52 & 94.07 & 81.98 & 76.06 & 94.61 & 97.02 & 90.75 & 98.05 & 94.90 &    
  85.54 & 80.36 & 68.37 & 86.17 \\                                                      
  Order~II & 50.02 & 92.57 & 85.49 & 77.18 & 95.24 & 96.88 & 90.58 & 96.16 & 93.62 &    
  86.92 & 78.51 & 67.97 & 85.74 \\                                                      
  \bottomrule     
  \end{tabular}%                                                                        
  }               
  \end{table}
   \section{Exemplar Allocation Formula}
  \label{app:allocation}

  We allocate the fixed buffer budget $B{=}11{,}000$ across seen tasks
  proportionally to class count: $b_t = \lfloor B \cdot C_t / \sum_{j \le t} C_j
  \rfloor$, where $C_t$ is the number of classes in task $t$. Existing tasks are
  downsampled to their updated $b_t$ via random sampling whenever a new task is
  added. At the end of all 11 tasks the total class count is $\sum C_j = 1{,}201$,
  yielding approximately $9.2$ exemplars per class across the buffer
  (\Cref{tab:iterations}). We tested uniform allocation ($b_t = B/T$) as well;
  both schemes match within $0.01$ on Last and Transfer (\Cref{tab:ablation},
  rows 2--3), so the choice is not load-bearing for the reported results.


  \section{Per-Task Training Iterations and Exemplar Allocation}
  \label{app:iterations}

  \begin{table}[h]
  \centering
  \caption{Per-task iteration counts and exemplar buffer allocation for all
  full-method experiments. Iterations were set to approximate CE loss
  convergence and fixed before the 11-task run. Buffer allocation follows
  $b_t = \lfloor B \cdot C_t / \sum_{j} C_j \rfloor$ with $B{=}11{,}000$
  and total classes $\sum C_j = 1{,}201$, yielding approximately $9.2$
  exemplars per class for all tasks.}
  \label{tab:iterations}
  \begin{tabular}{lcccc}
  \toprule
  \textbf{Task} & \textbf{Classes} & \textbf{Iterations} & \textbf{Buffer} & \textbf{Ex/class} \\
  \midrule
  FGVC-Aircraft    & 100 & 3{,}000 & 916   & 9.2 \\
  Caltech-101      & 101 & 1{,}000 & 925   & 9.2 \\
  CIFAR-100        & 100 & 1{,}500 & 916   & 9.2 \\
  DTD              &  47 & 1{,}500 & 430   & 9.1 \\
  EuroSAT          &  10 & 1{,}000 &  91   & 9.1 \\
  Flowers-102      & 102 & 1{,}500 & 934   & 9.2 \\
  Food-101         & 101 & 1{,}500 & 925   & 9.2 \\
  MNIST            &  10 &     800 &  91   & 9.1 \\
  Oxford-IIIT Pet  &  37 & 1{,}500 & 339   & 9.2 \\
  StanfordCars     & 196 & 3{,}000 & 1{,}795 & 9.2 \\
  SUN397           & 397 & 5{,}000 & 3{,}638 & 9.2 \\
  \midrule
  \textbf{Total}   & \textbf{1{,}201} & & \textbf{11{,}000} & \\
  \bottomrule
  \end{tabular}
  \end{table}

 
  \input{checklist}

  \end{document}
